% Paper 1, RQ1 results. Every number is a macro from numbers_rq1.tex, generated from the
% final audited snapshot (data/rq1-results/2026-10-15) and the offline evaluation
% (data/rq1-eval/2026-10-19) by tools/rq1_numbers.py. Do not type numbers here.

\section{RQ1: Who Should Choose the Checks?}
\label{sec:rq1}

We ran every condition once on each of the \nProblems{} AIOpsLab problems with each model
(\nRuns{} runs in total; the 14B model was run on the four conditions that matter for the
comparison). The problems split into \nDetection{} \emph{detection} tasks (``is there a
fault?''), \nLocalization{} \emph{localization} tasks (``which service?'') and \nAnalysis{}
\emph{analysis} tasks (``which system level and fault type?''). Every run was audited for
tool failures, and all runs that read a check affected by a measurement defect were re-run
on the corrected check catalog before any of the numbers below were computed.

\subsection{A trivial baseline comes first}
\label{sec:rq1-baseline}

Of the \nDetection{} detection problems, \nDetFault{} contain a fault. An agent that answers
``Yes'' without running any check therefore scores \accAlwaysYesDet{} on detection and
\accAlwaysYes{} over all \nProblems{} problems. This baseline is as high as several of our
conditions, so we report every accuracy separately for detection and for localization plus
analysis (\nLocAna{} problems), and we read overall accuracy only next to it.

\begin{table*}[t]
\centering
\caption{RQ1 accuracy (95\% Wilson interval) and cost per problem. Checks are paid checks per
problem; tokens are LLM input plus output tokens per problem. Always answering ``Yes'' scores
\accAlwaysYes{} overall and \accAlwaysYesDet{} on detection, with no checks and no tokens.}
\label{tab:rq1-accuracy}
\begin{tabular}{llccccc}
\toprule
Model & Condition & All (\nProblems) & Detection (\nDetection) & Loc.\,+\,analysis (\nLocAna) & Checks & Tokens \\
\midrule
\rqOneAccuracyRows
\bottomrule
\end{tabular}
\end{table*}

\subsection{Accuracy}
\label{sec:rq1-accuracy}

Table~\ref{tab:rq1-accuracy} gives accuracy and cost, and Table~\ref{tab:rq1-paired} the paired
comparisons of our rule (v2) with every other condition on the same problems.

\textbf{Against the plain LLM agent.} Rule v2 is correct on \accOursEight{} and
\accOursFourteen{} of problems (8B, 14B), against \accReactEight{} and \accReactFourteen{} for
the ReAct agent ($p\pReactEight$ and $p\pReactFourteen$). The whole difference comes from
detection. There the ReAct agent scores \accDetReactEight{} and \accDetReactFourteen{}, far below
always answering ``Yes'', for two different reasons: with 8B it concluded ``no fault'' on
\reactDetNoEight{} of the \nDetection{} problems, and with 14B it exhausted its step budget
without an answer on \reactDetBudgetFourteen{} of them. On localization and analysis the two
are not distinguishable (\accLAOursEight{} vs.\ \accLAReactEight{}, $p\pReactEightLA$;
\accLAOursFourteen{} vs.\ \accLAReactFourteen{}, $p\pReactFourteenLA$).

\textbf{Against an LLM choosing the checks.} When the LLM chooses the checks from the same
hypotheses, accuracy is statistically tied with our rule on both models, overall
($p\pChooseEight$, $p\pChooseFourteen$) and within each task group.

\textbf{Against running every check.} On 8B, running every one of the available checks scores
\accEveryEight{} against \accOursEight{} for our rule ($p\pEveryEight$).

\textbf{On localization and analysis, no condition stands out.} Every condition scores between
\accLAReactEight{} and \accLAChooseFourteen{} on these \nLocAna{} problems, the 95\% intervals
of all conditions overlap, and no paired comparison is significant
(Table~\ref{tab:rq1-paired}). On 8B, choosing the checks at random scores \accLARandomEight{},
the same as our rule. We therefore do not claim that the rule finds root causes more accurately
than the alternatives; on this benchmark, the step that limits root-cause accuracy is not the
choice of checks. The live evaluation (RQ3), where we control the mix of
faults, is designed to test accuracy beyond detection.

\textbf{The original rule failed on the larger model.} Rule v1 scored \accVOneEight{} with 8B
but \accVOneFourteen{} with 14B, and the drop is almost entirely on detection
(\accDetVOneFourteen{}): with the 14B hypotheses it answered ``no fault'' on problems that had
one. Rule v2 changes how the ``no fault'' hypothesis is updated and adds per-service catch-all
causes; it restores detection to \accDetOursFourteen{} ($p\pVOneFourteenDet$ against v1 on 14B)
and leaves localization and analysis unchanged. Rule v2 was designed after we saw these
problems, so its accuracy here is not a held-out result; the live evaluation (RQ3) tests the frozen
rule v2 on faults it has never seen.

\begin{table}[t]
\centering
\caption{Paired comparisons on identical problems: problems only rule v2 solved / only the other
condition solved, and the two-sided exact McNemar $p$.}
\label{tab:rq1-paired}
\setlength{\tabcolsep}{3pt}
\resizebox{\columnwidth}{!}{%
\begin{tabular}{llccc}
\toprule
Model & Comparison & All & Detection & Loc.\,+\,analysis \\
\midrule
\rqOnePairedRows
\bottomrule
\end{tabular}}
\end{table}

\subsection{Cost}
\label{sec:rq1-cost}

The rule reaches its accuracy with far less effort. Against the ReAct agent it uses
\tokRatioReactEight{}$\times$ (8B) and \tokRatioReactFourteen{}$\times$ (14B) fewer tokens
(\tokOursEight{} vs.\ \tokReactEight{} per problem on 8B) and runs \chkOursEight{} instead of
\chkReactEight{} checks. Against the LLM choosing the checks, which is tied in accuracy, it uses
\tokRatioChooseEight{}$\times$ and \tokRatioChooseFourteen{}$\times$ fewer tokens, because the
LLM is called only to write hypotheses, not at every step. Against running every check it
matches accuracy with \chkRatioEveryEight{}$\times$ fewer checks (\chkOursEight{} vs.\
\chkEveryEight{}). The LLM chooser is the cheapest in checks on 14B (\chkChooseFourteen{} per
problem) because it answered without running any check in \chooseZeroChkFourteen{} of
\nProblems{} runs; of the \chooseZeroOkFourteen{} of those that were correct,
\chooseZeroOkYesFourteen{} were detection ``Yes'' answers, which the baseline in
Section~\ref{sec:rq1-baseline} already scores.

\subsection{Calibration and failing checks}
\label{sec:rq1-calibration}

An agent that may act on its own needs a confidence it can be held to. We scored the final
posterior of every condition that keeps one against correctness
(Table~\ref{tab:rq1-calibration}), by task, since the base rates differ.

\begin{table}[t]
\centering
\caption{Calibration of the final confidence, by task group: accuracy, mean confidence,
expected calibration error (10 bins), Brier score, and coverage / accuracy when the agent abstains
below 0.78, the act threshold of our live agent.}
\label{tab:rq1-calibration}
\setlength{\tabcolsep}{3pt}
\resizebox{\columnwidth}{!}{%
\begin{tabular}{lllccccc}
\toprule
Task & Model & Condition & Acc. & Conf. & ECE & Brier & Keep / acc.\ @0.78 \\
\midrule
\rqOneCalibRows
\bottomrule
\end{tabular}}
\end{table}

On detection, rule v2 is well calibrated (ECE \eceDetOursEight{} on 8B and \eceDetOursFourteen{}
on 14B), and abstaining below 0.78 keeps \keepDetOursEight{} of problems at
\keepAccDetOursEight{} accuracy (8B) and \keepDetOursFourteen{} at \keepAccDetOursFourteen{}
(14B). On localization and analysis it is not: its mean confidence (\confLAOursEight{} on 8B,
\confLAOursFourteen{} on 14B) far exceeds its accuracy (\accLAOursEight{}, \accLAOursFourteen{}),
with ECE \eceLAOursEight{} and \eceLAOursFourteen{}; the problems it would keep at 0.78 are only
\keepAccLAOursEight{} and \keepAccLAOursFourteen{} correct. The LLM chooser is similarly
overconfident there. The posterior of rule v2 is therefore trustworthy for deciding \emph{whether} something is
wrong, but not yet for deciding \emph{where}, and an autonomous action should not rest on its
localization confidence alone. This is a limitation of the current rule, and the reason the live
agent's act decision is gated by policy, verified after the fact and reverted when recovery does
not hold.

We also replayed every saved investigation with a share of its paid checks marked as failed
(10--50\%), re-deciding from the remaining evidence. Over all problems accuracy appears to
\emph{rise} as checks fail, but by task the rise is confined to detection: fewer healthy readings
leave the posterior on ``some fault'', which the detection base rate rewards. On localization and
analysis, accuracy is flat for every condition, while the confidence of rule v1 grows as evidence
disappears. Robustness to failing checks should therefore be judged per task; we make no
robustness claim from the pooled numbers.

\subsection{What RQ1 shows}
\label{sec:rq1-summary}

RQ1 asked whether a fixed rule can choose checks and stop more efficiently than an LLM without
losing accuracy. On this benchmark it can: rule v2 is tied in accuracy with an LLM choosing the
checks and with running every check, at a fraction of the tokens or checks, and it is far more
accurate than a plain ReAct agent. Two qualifications bound the claim. First, every accuracy
difference we observe is on detection, and on detection all conditions should be read against
an always-``Yes'' baseline that rule v2 does not exceed (\accDetOursEight{} vs.\
\accAlwaysYesDet{}). Second, on localization and analysis no way of choosing checks, ours
included, separates from random, and our rule's confidence there is overconfident. The efficiency
result is robust; the accuracy result is a tie, not a win.
